\chapter{คู่มือปฏิบัติการวัดละเอียดทางฟิสิกส์}
\label{app:lab_manual}
\index{คู่มือปฏิบัติการวัดละเอียด}

คู่มือปฏิบัติการฉบับนี้จัดทำขึ้นเพื่อฝึกทักษะการใช้เครื่องมือวัดละเอียดพื้นฐานสำหรับนักศึกษาสาขาวิทยาศาสตร์ชีวภาพ จุลชีววิทยา และสิ่งแวดล้อม ประกอบด้วยเวอร์เนียร์คาลิปเปอร์ ไมโครมิเตอร์ และสเฟียโรมิเตอร์

\section{ปฏิบัติการที่ 1: การวัดขนาดวัตถุอย่างละเอียด}

\subsection{วัตถุประสงค์การทดลอง}
\begin{enumerate}
    \item เพื่อศึกษาและฝึกทักษะการวัดขนาดวัตถุด้วยเวอร์เนียร์คาลิปเปอร์ (Vernier Caliper) และไมโครมิเตอร์ (Micrometer)
    \item เพื่อศึกษาการใช้สเฟียโรมิเตอร์ (Spherometer) ในการวัดรัศมีความโค้งของเลนส์และผิวโค้งทรงกลม
    \item เพื่อฝึกการวิเคราะห์และรายงานค่าความคลาดเคลื่อนตามหลักเลขนัยสำคัญ
\end{enumerate}

\subsection{ทฤษฎีและหลักการของเครื่องมือวัด}
\begin{enumerate}
    \item \textbf{เวอร์เนียร์คาลิปเปอร์:} เครื่องมือวัดความยาว ความหนา เส้นผ่านศูนย์กลางภายใน ภายนอก และความลึก มีความละเอียดของสเกลเวอร์เนียร์ (Least Count) คำนวณจาก:
    \begin{equation}
        LC = \frac{\text{ความยาว 1 ช่องสเกลหลัก}}{\text{จำนวนช่องสเกลเวอร์เนียร์}} = \frac{1\,\text{mm}}{20} = 0.05\,\text{mm}
    \end{equation}
    
    \item \textbf{ไมโครมิเตอร์:} เครื่องมือวัดละเอียดที่อาศัยหลักการเกลียวละเอียด (Screw Pitch) หนึ่งรอบการหมุนเลื่อนแกนวัดได้ $0.5\,\text{mm}$ ปลอกหมุนแบ่งเป็น 50 ช่องสเกล ความละเอียดต่ำสุดคือ:
    \begin{equation}
        LC = \frac{0.5\,\text{mm}}{50} = 0.01\,\text{mm} = 10\,\mu\text{m}
    \end{equation}
\end{enumerate}

\begin{table}[H]
\centering
\small
\begin{tabularx}{\textwidth}{l c c c Y}
\toprule
\textbf{เครื่องมือวัด} & \textbf{สเกลหลัก} & \textbf{สเกลละเอียด} & \textbf{Least Count} & \textbf{การนำไปใช้ในงานชีวภาพ} \\
\midrule
ไม้บรรทัดเหล็ก & 1 mm & - & 0.5 mm & วัดความยาวตัวอย่างราก/ลำต้น \\
เวอร์เนียร์คาลิปเปอร์ & 1 mm & 20 หรือ 50 ช่อง & 0.05 หรือ 0.02 mm & วัดขนาดเส้นผ่านศูนย์กลางผล/หลอดทดลอง \\
ไมโครมิเตอร์ & 0.5 mm & 50 ช่อง & 0.01 mm & วัดความหนาของแผ่นสไลด์/ใบพืช \\
สเฟียโรมิเตอร์ & 1 mm & 100 ช่อง & 0.005 mm & วัดรัศมีความโค้งผิวเลนส์กล้องจุลทรรศน์ \\
\bottomrule
\end{tabularx}
\caption{คุณลักษณะและความละเอียดของเครื่องมือวัดละเอียดในห้องปฏิบัติการ}
\end{table}

\subsection{ตารางบันทึกผลการทดลอง}

\begin{table}[H]
\centering
\small
\begin{tabularx}{\textwidth}{c Y c c c}
\toprule
\textbf{ครั้งที่} & \textbf{ชิ้นงานตัวอย่าง} & \textbf{ค่าที่วัด (mm)} & \textbf{$\bar{x}$ (mm)} & \textbf{$S_D$ (mm)} \\
\midrule
1 & ความหนากระจกสไลด์ (ไมโครมิเตอร์) & 1.22 & \multirow{3}{*}{1.223} & \multirow{3}{*}{$\pm 0.006$} \\
2 & ความหนากระจกสไลด์ (ไมโครมิเตอร์) & 1.23 & & \\
3 & ความหนากระจกสไลด์ (ไมโครมิเตอร์) & 1.22 & & \\
\midrule
1 & เส้นผ่านศูนย์กลางหลอดทดลอง (เวอร์เนียร์) & 15.45 & \multirow{3}{*}{15.47} & \multirow{3}{*}{$\pm 0.020$} \\
2 & เส้นผ่านศูนย์กลางหลอดทดลอง (เวอร์เนียร์) & 15.50 & & \\
3 & เส้นผ่านศูนย์กลางหลอดทดลอง (เวอร์เนียร์) & 15.45 & & \\
\bottomrule
\end{tabularx}
\caption{ตัวอย่างตารางบันทึกผลการวัดละเอียดและการคำนวณค่าทางสถิติ}
\end{table}
